% ID: SEL-D0079-OB-E04
% Tanda S079.
\textbf{a)} Contrastamos $H_0:p\ge0{,}26$ frente a $H_1:p<0{,}26$. La alternativa indica un contraste unilateral a la izquierda. Usamos la aproximación normal para una proporción, suponiendo una muestra aleatoria de usuarios independientes.

En la frontera de la hipótesis nula, $p_0=0{,}26$, el estadístico es
\[
 Z=\frac{\hat p-0{,}26}{\sqrt{0{,}26\cdot0{,}74/1000}}.
\]
Bajo $p=p_0$ se aproxima por una normal estándar. Con significación $\alpha=0{,}05$, la región crítica es
\[
 \boxed{Z<-1{,}64485},
 \qquad\text{equivalentemente}\qquad\hat p<0{,}23718.
\]
Para $p>p_0$, el rechazo por proporciones bajas es menos probable, por lo que se usa la frontera para fijar el contraste.

\textbf{b)} La proporción observada es $\hat p=240/1000=0{,}24$. Sustituyendo,
\[
 Z_{\mathrm{obs}}\approx-1{,}4419>-1{,}64485.
\]
El resultado queda fuera de la región crítica y no se rechaza $H_0$ al $5\%$. Los datos son compatibles con el informe del Ayuntamiento a ese nivel. Esto no prueba que la proporción sea al menos el $26\%$. El valor $p$ aproximado es $\Phi(-1{,}4419)\approx0{,}0747$.
